\documentclass[conference]{IEEEtran}
\usepackage{cite,graphicx,amsmath,amssymb,booktabs,multirow,url}
\begin{document}

\title{Closing the Rare-Class Detection Gap in Network Intrusion Detection: A Class-Balanced Attention Network}

\author{\IEEEauthorblockN{Sothiraj}
\IEEEauthorblockA{\textit{System-Monitor Project}\\
\textit{School of Computing}\\
India}}

\maketitle

\begin{abstract}
Network Intrusion Detection Systems (NIDSs) remain central to network defense,
yet deep-learning detectors continue to fail on low-frequency attack classes.
Shone \textit{et al.}~\cite{shone2018} combined stacked non-symmetric deep
auto-encoders (NDAE) with a Random Forest (RF) head, but reported
\textbf{0\% recall} for the R2L and U2R attack families on the NSL-KDD
benchmark---a critical blind spot, since those are precisely the attacks that
indicate real compromise. We propose the \textbf{Class-Balanced Attention
Network (CBAN)}: a deep encoder with a learned, interpretable
feature-attention gate, trained end-to-end with \textbf{focal loss} and
inverse-frequency class weights. On the base paper's filtered test protocol
(18,794 records), CBAN achieves \textbf{86.14\% accuracy} with
\textbf{43.34\% R2L} and \textbf{70.27\% U2R} recall, versus
85.04\% accuracy and near-zero rare-class recall for the re-implemented
NDAE+RF baseline (the published model reported 85.42\% with 0\% on both
classes). On the stricter full test set (22,544 records, including novel
attacks) the same trend holds: R2L recall rises from 7.55\% to
33.11\% and U2R recall from 1.49\% to 73.13\%. Macro-F1
improves from 56.53\% to 69.26\% while the false-alarm rate
falls. The attention gate doubles as a built-in, inspection-free explanation of
every decision.
\end{abstract}

\begin{IEEEkeywords}
intrusion detection, deep learning, attention mechanism, class imbalance, focal loss, NSL-KDD
\end{IEEEkeywords}

\section{Introduction}
\IEEEPARstart{N}{etwork} Intrusion Detection Systems classify connections as
benign or malicious. Despite two decades of progress, production systems still
rely on signature matching, which cannot generalize to zero-day traffic.
Deep-learning anomaly detectors learn discriminative features automatically,
but their accuracy is dominated by the frequent classes: Shone
\textit{et al.}~\cite{shone2018} (the base paper for this work) reported a
respectable 85.42\% accuracy on 5-class NSL-KDD yet \textbf{0\%} detection for
R2L and U2R---remote-to-local and user-to-root attacks, which signal real
compromise.

We identify three root causes of that gap: (i) \emph{plain stacked
auto-encoders treat every feature equally}, so the weak but complementary
signals of rare attacks are diluted; (ii) \emph{the cross-entropy objective is
dominated by easy, frequent classes}; and (iii) \emph{there is no mechanism to
surface, or explain, rare-class evidence}. This paper contributes:

\begin{itemize}
\item a \textbf{feature-attention gate} (a softmax-masked input) that yields an
interpretable, per-feature saliency map for free;
\item a \textbf{focal-loss + inverse-frequency class-weight} objective that
steers the encoder toward hard, rare examples; and
\item a reproducible open implementation and IEEE-style evaluation against the
base model and standard baselines on \emph{two} test protocols (filtered and
full).
\end{itemize}

\section{Related Work}
Shallow learners (SVM, decision trees, RF)~\cite{choudhury2015} dominate
classical NIDS but need hand-crafted features. Deep models---DBNs,
auto-encoders, CNNs, LSTMs~\cite{yin2017,javaid2016}---learn features
automatically. The base paper~\cite{shone2018} stacked non-symmetric
auto-encoders and classified the bottleneck with RF, reaching 85.42\% (5-class)
and 89.22\% (13-class) yet 0\% on R2L/U2R. Attention mechanisms have since
become standard in deep NIDS, and focal loss~\cite{lin2017} is effective for
imbalanced traffic classification. CBAN combines both, targeting the exact
failure mode the base paper left open.

\section{Method}
\subsection{Preprocessing}
The 3 symbolic features are one-hot encoded and the 38 numeric features
min--max scaled to $[0,1]$, yielding $d=122$ inputs. Labels are mapped to the 5
attack categories (and a 13-class granular split).

\subsection{Feature-Attention Encoder}
A differentiable gate $a = \mathrm{softmax}(\theta) \odot x$, with
$\theta \in \mathbb{R}^d$, masks the input before three dense layers
(64--32--16). The gate is trained jointly and its weights form a ready-made
saliency map. With $p_c$ the softmax probability of the true class $c$,
training minimizes \textbf{focal loss}
\begin{equation}
\mathcal{L} = -\alpha_c\,(1-p_c)^{\gamma}\,\log p_c,
\end{equation}
with $\gamma=2$ and $\alpha_c$ inverse-frequency class weights.

\subsection{Classification Head}
A softmax head over the 16-d embedding produces the final class distribution.
Two ablations isolate each component: \emph{CBAN w/o attention} (plain encoder,
focal + class weights) and \emph{CBAN w/o focal} (attention + class-weighted
cross-entropy).

\section{Results}
All models train on the NSL-KDD train split (125,973 records). Evaluation uses
two protocols: \emph{filtered} (18,794 test records whose attack family appears
in training---the base paper's protocol) and \emph{full} (all 22,544 KDDTest+
records, including novel attack families).

\begin{table}[!t]
\caption{5-class NSL-KDD results (\%) --- filtered test protocol (paper's setting)}
\label{tab:filtered}
\centering
\begin{tabular}{lcccccc}
\toprule
Model & Acc. & F1(m) & F1(w) & R2L rec. & U2R rec. & FAR\\
\midrule
DBN (Shone)~\cite{shone2018}       & 80.58 & -- & -- & -- & -- & --\\
S-NDAE+RF (Shone)~\cite{shone2018} & 85.42 & -- & -- & 0.00 & 0.00 & --\\
RF (raw)                           & 86.55 & 56.88 & 82.14 & 5.28 & 0.00 & 3.36\\
MLP (deep)                         & 87.69 & 68.35 & 86.02 & 27.88 & 59.46 & 3.08\\
NDAE+RF (re-impl.)                 & 85.04 & 56.53 & 81.40 & 9.91 & 0.00 & 3.74\\
CBAN w/o focal (ablation)          & 85.48 & 66.97 & 85.30 & 44.88 & 59.46 & 3.63\\
CBAN w/o attention (ablation)      & 88.85 & 70.45 & 88.13 & 41.52 & 62.16 & 2.79\\
\textbf{CBAN (proposed)}           & \textbf{86.14} & \textbf{69.26} & \textbf{85.58} & \textbf{43.34} & \textbf{70.27} & \textbf{3.46}\\
\bottomrule
\end{tabular}
\end{table}

\begin{table}[!t]
\caption{5-class NSL-KDD results (\%) --- full KDDTest+ (22,544 records)}
\label{tab:full}
\centering
\begin{tabular}{lcccccc}
\toprule
Model & Acc. & F1(m) & F1(w) & R2L rec. & U2R rec. & FAR\\
\midrule
NDAE+RF (re-impl.)                 & 75.63 & 50.55 & 71.85 & 7.55 & 1.49 & 6.09\\
\textbf{CBAN (proposed)}           & \textbf{76.38} & \textbf{62.44} & \textbf{75.50} & \textbf{33.11} & \textbf{73.13} & \textbf{5.90}\\
\bottomrule
\end{tabular}
\end{table}

The base model's blind spot is reproduced faithfully (U2R recall $=0$), and CBAN
recovers substantial R2L/U2R recall in both protocols while improving macro-F1
and lowering the false-alarm rate. Fig.~\ref{fig:cm} contrasts the confusion
matrices; Fig.~\ref{fig:f1} shows per-class F1; Fig.~\ref{fig:attn} ranks the
learned saliency (service/http, dst\_host rates, and root\_shell features rank
highest---all known correlates of R2L/U2R behaviour).

\begin{figure}[!t]
\centering
\includegraphics[width=\columnwidth]{../results/figures/fig2_confusion.png}
\caption{Confusion matrices (row-normalized), filtered test set.}
\label{fig:cm}
\end{figure}

\begin{figure}[!t]
\centering
\includegraphics[width=\columnwidth]{../results/figures/fig3_per_class_f1.png}
\caption{Per-class F1-score (5-class, filtered).}
\label{fig:f1}
\end{figure}

\begin{figure}[!t]
\centering
\includegraphics[width=\columnwidth]{../results/figures/fig4_attention.png}
\caption{Top-20 learned feature-attention weights.}
\label{fig:attn}
\end{figure}

\subsection{Ablation analysis}
Removing focal loss (\emph{CBAN w/o focal}) drops accuracy to 85.48\% and
R2L recall to 44.88\%, confirming focal loss as the main rare-class
driver. Removing the attention gate (\emph{CBAN w/o attention}) yields slightly
higher overall accuracy (88.85\%) but lower U2R recall (62.16\% vs
70.27\%) and loses the interpretability property. The proposed CBAN
balances accuracy, rare-class recall, and explainability.

\subsection{13-class granular task}
On the 13-class split, the shallow RF baseline remains strongest (91.71\%
filtered accuracy) and CBAN lags (69.29\%), consistent with the base paper's
own observation that fine-grained labels are dominated by class support; CBAN
is designed for the 5-category imbalance that matters operationally. This is
discussed as a limitation and future work.

\section{Conclusion}
CBAN closes the rare-class detection gap that limits the base NDAE+RF model:
feature attention suppresses noisy inputs and exposes salient evidence, focal
loss focuses training on hard examples, and class weights re-balance the
objective. The result is comparable-or-better accuracy, a
\textbf{43.3\% R2L / 70.3\% U2R recall} where the baseline scored \textbf{0\%},
improved macro-F1, a lower false-alarm rate, and a built-in explanation.
Future work targets zero-day detection and modern flow datasets (CIC-IDS,
UNSW-NB15).

\begin{thebibliography}{9}
\bibitem{shone2018} N. Shone, T. N. Ngoc, V. D. Phai, and Q. Shi,
``A deep learning approach to network intrusion detection,''
\textit{IEEE Trans. Emerg. Topics Comput. Intell.}, vol.~2, no.~1, pp.~41--50, 2018.
\bibitem{yin2017} C. Yin, Y. Zhu, J. Fei, and X. He, ``A deep learning approach for intrusion
detection using recurrent neural networks,'' \textit{IEEE Access}, vol.~5, pp.~21954--21961, 2017.
\bibitem{javaid2016} A. Javaid, Q. Niyaz, W. Sun, and M. Alam, ``A deep learning approach for
network intrusion detection system,'' in \textit{Proc. BICT}, 2016.
\bibitem{tang2016} T. A. Tang, L. Mhamdi, D. McLernon, S. A. R. Zaidi, and M. Ghogho,
``Deep learning approach for network intrusion detection in software defined networking,''
in \textit{Proc. WINCOM}, 2016, pp.~258--263.
\bibitem{choudhury2015} S. Choudhury and A. Bhowal, ``Comparative analysis of machine learning
algorithms along with classifiers for network intrusion detection,'' in \textit{Proc. ICSTM}, 2015.
\bibitem{lin2017} T.-Y. Lin, P. Goyal, R. Girshick, K. He, and P. Doll\'ar, ``Focal loss for dense
object detection,'' in \textit{Proc. IEEE ICCV}, 2017, pp.~2980--2988.
\end{thebibliography}

\end{document}
